\section{Minimum subspace trades}
\label{sec:circuits}

A point trade in \(\cB_i\) is an integral relation written with disjoint
positive and negative supports,
\[
  T=T^+-T^-=\sum_{B\in T^+}B-\sum_{B\in T^-}B
  \in K_i^{\ZZ}.
\]
Its volume is \(|T^+|=|T^-|\), and its total support is
\(|T^+|+|T^-|\).  This distinction matters here: the minimum
\(T_2(1,4,8)\) trades have volume \(3\) and total support \(6\).

\subsection{The opposite-regulus model}

Choose vector subspaces \(Z<U<\FF_2^8\) with
\(\dim Z=2\) and \(\dim U=6\).  A four-space \(B\) satisfying
\(Z<B<U\) projects to a projective line \(B/Z\) in
\(\operatorname{PG}(U/Z)\cong\operatorname{PG}(3,2)\).
Three pairwise skew lines form one regulus of the hyperbolic quadric
\(Q^+(3,2)\); the three lines meeting all of them form its opposite
regulus.  Lifting the two reguli gives two triples of four-spaces
\[
  \{B_1,B_2,B_3\},
  \qquad
  \{B'_1,B'_2,B'_3\},
\]
with the same point-incidence sum.

\begin{lemma}[Primitive regulus circuit]
\label{lem:regulus-circuit}
Let \(N\) be the \(255\times6\) point-incidence matrix of the six lifted
blocks, ordered by the two reguli.  Then
\[
 \rank_{\QQ}N=5,\qquad
 \ker_{\ZZ}N
  =\ZZ(1,1,1,-1,-1,-1).
\]
Every five columns are independent, and the nonzero Smith factors of
\(N\) are \(1,1,1,1,1\).  Hence the displayed trade is a primitive
unimodular matroid circuit and a Graver element.
\end{lemma}

\begin{proof}
In the quotient \(\operatorname{PG}(3,2)\), each point of the hyperbolic
quadric lies on one line from each regulus, giving the displayed kernel
vector.  Direct row reduction of the \(15\)-point quotient incidence
matrix has rank five, and deleting any one line leaves five independent
columns.  Lifting through \(Z<U\) repeats rows but introduces no new
column relation.  The active Smith form is
\(\operatorname{diag}(1,1,1,1,1)\), so the one-dimensional integral
kernel is primitive and generated by the sign vector above.  A primitive
support-minimal relation is a Graver element.
\end{proof}

Krotov's minimum-trade theorem applies with
\[
 0\le t=1<k=4<v-t=7.
\]
Its support bound specializes to
\(\prod_{j=0}^{1}(1+2^j)=6\), hence volume three, and its \(t=1\)
classification identifies every trade attaining this bound with the
opposite-regulus construction above \cite{Krotov2017}.  We use that
classification for exhaustiveness, while \cref{lem:regulus-circuit}
supplies the integral circuit statement.

\subsection{Exact census}

Let \(\cC_i\) be the set of unoriented minimum trades wholly contained in
\(\cB_i\).  Exhaustive enumeration first chooses \((Z,U)\), constructs
the two reguli in \(U/Z\), and then tests membership of all six lifts in
\(\cB_i\).  Canonical sorting removes repeated presentations, and the
common group then partitions the result into orbits.

\begin{proposition}[Minimum-trade census]
\label{prop:circuit-census}
The numbers of minimum trades and their \(\GammaG\)-orbit counts are
\[
\begin{array}{c|rrr}
 &\cB_1&\cB_2&\cB_3\\
\midrule
|\cC_i|&
2\,271\,676&2\,213\,026&2\,216\,868\\
|\cC_i/\GammaG|&
11\,274&11\,001&11\,033.
\end{array}
\]
Every trade in the census has coefficient gcd one and active Smith form
\(1^5\).  Thus every local arithmetic multiplicity is one.
\end{proposition}

\begin{proof}
Krotov's \(t=1\) classification reduces the search to the
\((Z,U)\)-regulus templates, so the enumeration is exhaustive rather
than a bounded-support sample.  Three independent passes---template
generation, direct six-block membership, and orbit closure---give the
displayed totals.  For each accepted template,
\cref{lem:regulus-circuit} gives the coefficient gcd and Smith form.
The frozen counts, orbit representatives, and replay hashes form the
census certificate of \cref{app:certificates}.
\end{proof}

\begin{remark}
The configurations are separated by how many minimum circuits they
contain and how those circuits are organized into group orbits.  They are
not separated by a different six-block local arithmetic type.
Nothing in this section classifies minimum \(T_2(2,4,8)\) trades; for that
parameter the same support formula gives \(30\), or volume \(15\).
\end{remark}
